\section{国内外本学科领域的发展现状与趋势}

本章围绕课题的三个研究点，分别梳理相关领域的研究背景、研究现状与趋势、以及尚存的关键挑战。其中，2.1 节对应语义漏洞挖掘研究点，2.2 节与 2.3 节分别对应拟开展的漏洞自动修复与漏洞修复验证研究。

\subsection{基于智能体的语义漏洞挖掘}

\subsubsection{研究背景}

漏洞挖掘旨在从目标软件中自动发现未知缺陷。经过数十年发展，以模糊测试（Fuzzing）、符号执行（Symbolic Execution）、污点分析（Taint Analysis）为代表的技术已在内存破坏类漏洞（缓冲区溢出、释放后使用等）上取得显著成功。然而，\textbf{在语义/逻辑漏洞（semantic / logic vulnerability）方面，自动化挖掘能力仍然薄弱}。

语义漏洞指程序在内存安全层面"合法"、但在设计意图或业务逻辑层面"错误"的缺陷，其典型类别包括：认证与授权绕过、状态机错误、配置驱动缺陷、第三方组件使用违规、跨组件逻辑错误等。与内存破坏漏洞不同，语义漏洞\textbf{不产生崩溃信号}，无法被 AddressSanitizer 等通用 oracle 直接捕获；其判定往往需要人工或半自动构造的策略/规约（policy/specification），这是语义漏洞挖掘区别于内存破坏挖掘的根本难点。

在 IoT 与嵌入式固件场景下，该问题尤为突出。固件通常包含大量由\textbf{内部状态、外设寄存器、NVRAM 配置、时序事件}共同决定的程序路径，这些路径对传统的输入驱动 fuzzing 而言几乎不可达。我们将这类"不受外部输入直接控制、由系统内部状态/配置/时序决定的变量"统称为\textbf{自因变量（Self-Caused Variables）}。自因变量普遍存在于认证逻辑、状态机转移、外设交互与配置处理中，是语义漏洞挖掘的核心障碍。

与此同时，大语言模型与智能体展现出前所未有的自然语言规约理解、跨抽象层推理与代码合成能力，为突破语义漏洞挖掘瓶颈提供了新的可能。智能体可以阅读协议 RFC、数据手册、配置说明等自然语言文档，自动提取安全不变量作为 oracle；可以在 fuzzing 循环中观察覆盖率停滞，主动推断需要翻转哪个内部状态变量，并合成对应的环境事件序列；还可以跨驱动层、RTOS 任务、应用层进行语义推理，识别跨层逻辑错误。

\subsubsection{研究现状与趋势}

\paragraph{（一）语义漏洞检测的传统方法}

围绕语义漏洞检测，已有研究沿多条技术路线展开：

\begin{itemize}
  \item \textbf{静态符号执行与模型检验}：Firmalice~\cite{firmalice} 以"认证绕过"为唯一语义，结合特权模型进行符号执行，是该方向的奠基性工作；FIE~\cite{fie} 面向 MSP430 固件进行全系统符号执行。此类方法受限于路径爆炸与环境建模难度，难以扩展到大型固件。
  \item \textbf{静态污点与数据流分析}：KARONTE~\cite{karonte} 通过跨二进制污点分析检测固件多进程通信中的逻辑缺陷；Operation Mango~\cite{operationmango}、FIRE~\cite{fire}、OctopusTaint~\cite{octopustaint}、UVScan~\cite{uvscan} 等工作进一步提升了可扩展性与检测精度；LATTE~\cite{latte} 首次将 LLM 引入二进制污点分析，显著降低了规则定制成本。
  \item \textbf{固件重托管与覆盖率 Fuzzing}：FIRM-AFL~\cite{firmafl}、P2IM~\cite{p2im}、Laelaps~\cite{laelaps}、Fuzzware~\cite{fuzzware}、Ember-IO~\cite{emberio} 等工作通过外设建模与精确 MMIO 处理，显著提升了固件 fuzzing 的覆盖率；Firmvenom~\cite{firmvenom} 引入外设语义感知 fuzzing，进一步提升效率。
  \item \textbf{状态感知 Fuzzing}：AFLNet~\cite{aflnet} 以协议状态机引导 fuzzing；StateAFL~\cite{stateafl} 通过内存快照推断协议状态；SyzTrust~\cite{syztrust} 面向可信执行环境，明确提出"隐藏状态变量缺乏语义反馈"这一与自因变量高度对应的概念，并通过状态感知 fuzzing 缓解深状态不可达问题。
  \item \textbf{配置驱动 Fuzzing}：CONFU~\cite{confu} 等工作将配置文件直接作为 fuzz 输入，绕开"配置变量在运行期不变"带来的自因性。
\end{itemize}

\paragraph{（二）自因变量的建模与求解}

针对自因变量带来的挑战，已有研究从不同角度提出了建模与求解思路，如表~\ref{tab:selfcaused} 所示。总体上，SOTA 已从不把自因变量当作噪声，而是将其\textbf{显式对象化}——或提升为 fuzzer 一等输入，或通过符号/学习方法恢复其生成模型。但\textbf{尚缺乏统一抽象来整合 state-driven / config-driven / timing-driven / peripheral-driven 等不同形态}。

\begin{table}[htbp]
  \centering
  \caption{自因变量建模的代表性方法路线}
  \label{tab:selfcaused}
  \small
  \begin{tabular}{p{3.4cm}p{4.2cm}p{6.6cm}}
    \toprule
    \textbf{方法路线} & \textbf{代表工作} & \textbf{核心思想} \\
    \midrule
    外设状态机学习 & P2IM, Laelaps, Fuzzware, Firmvenom & 将 MMIO 建模为状态机，符号执行 + 自动机学习推断状态转移 \\
    运行时状态推断 & StateAFL, SyzTrust & 内存快照聚类推断隐藏状态，将状态作为 fuzzer 一等公民 \\
    配置显式化 & CONFU & 将配置变量作为独立输入维度 \\
    反向数据流切片 & KARONTE, Operation Mango & 从 sink 反向切片，识别不受 source 控制的变量并单独求解 \\
    混合执行 & Laelaps, FirmHybridFuzz & 输入不可解时切换符号执行，将内部状态变量显式符号化 \\
    \bottomrule
  \end{tabular}
\end{table}

\paragraph{（三）智能体在漏洞挖掘中的初步探索}

2024 年以来，智能体在漏洞挖掘中的应用开始爆发：FirmAgent~\cite{firmagent} 将 fuzzing 与 LLM 智能体结合，由 fuzzing 提供污点源，LLM 智能体进行上下文感知分析与 PoC 生成，在 14 个真实固件上发现 182 个漏洞，精度 91\%，获得 17 个 CVE——这是目前该方向最强的公开结果；FirmPilot~\cite{firmpilot} 基于多智能体协作进行固件重托管，用 LLM 推断外设与环境语义；ChatAFL~\cite{chatafl} 利用 LLM 提取协议语法、丰富种子、突破覆盖率平台；Big Sleep~\cite{bigsleep} 首次公开展示 AI 智能体在真实开源项目中发现未知内存安全漏洞；AIxCC~\cite{aixcc} 系统验证了智能体在"发现—修补—验证"全链路中的能力。此外，BinReX~\cite{binrex} 提出的"语义检索 + 可验证推理"范式表明：为智能体提供检索增强的语义导航与机器可验证的证据链，相比单纯扩大模型规模或工具集更能提升二进制分析效果，这一范式对语义漏洞挖掘具有直接借鉴意义。

\paragraph{（四）研究趋势}

研究范式正从"LLM 作为分类器/生成器嵌入传统 pipeline"转向"智能体作为分析主体，自主规划分析路径"；研究对象从内存破坏漏洞扩展到语义/逻辑漏洞；系统架构从单智能体演进到规划—分析—验证分工的多智能体；评估对象从合成基准扩展到真实固件镜像与大型开源仓库。

\subsubsection{研究挑战}

\begin{enumerate}
  \item \textbf{自因变量缺乏统一形式化}：state-driven、config-driven、timing-driven、peripheral-driven 等不同形态的自因变量在文献中各自为战，缺乏类似污点分析中 lattice 的统一抽象，导致方法难以泛化。
  \item \textbf{语义 oracle 的自动构造}：逻辑漏洞的判定准则难以自动抽取，Firmalice 的认证绕过 oracle 使用十年仍是孤例。如何利用智能体从自然语言文档中自动提取安全不变量，是关键开放问题。
  \item \textbf{覆盖率反馈对语义深度的"失明"}：coverage-guided 机制对"已覆盖代码内的逻辑错误"无感知，需要语义感知的反馈信号来替代或增强传统边覆盖。
  \item \textbf{跨层语义鸿沟}：固件漏洞常跨越"驱动层 $\rightarrow$ RTOS 任务 $\rightarrow$ 应用层 $\rightarrow$ 云端/APP"，现有工具基本局限于单层，缺乏跨抽象层的语义整合能力。
  \item \textbf{基准缺失}：缺乏公认的"语义漏洞 + 自因变量"基准，导致不同方法之间难以公平比较。
\end{enumerate}

\subsection{基于智能体的漏洞自动修复}

\subsubsection{研究背景}

漏洞修复是软件安全维护的重要环节，其目标不仅是阻止已知攻击输入触发异常，还需要消除导致危险行为的程序缺陷，并维持正常使用所依赖的功能与安全约定。漏洞自动修复（Automated Vulnerability Repair，AVR）将漏洞分析、补丁生成与补丁验证组织为自动化过程，以降低安全维护对人工分析的依赖。与面向一般软件缺陷的自动程序修复（Automated Program Repair，APR）相比，AVR 更强调在攻击者可控输入、异常状态及安全敏感操作下恢复必要的安全条件。因此，修复目标不能仅用代码修改是否合理、程序能否编译或已知失败是否消失来描述，还应考察危险行为是否被有效约束，以及相关合法行为是否得到保留~\cite{sokavr}。

漏洞修复通常发生在不完备信息条件下。漏洞报告、检测器告警和概念验证程序（Proof of Concept，PoC）能够提供分析线索，但不一定直接指出应修改的程序位置，也不一定完整描述正确修复应满足的条件。已有 APR 研究表明，通过用于生成补丁的测试集，并不意味着补丁能够推广到未见输入~\cite{smith2015cure}。在漏洞修复中，这一问题进一步表现为：已知 PoC 不再触发，可能只是因为某个特定输入或执行分支被阻断，而非漏洞根因已经消除。由此，AVR 实际上包含两个相互关联、但不能相互替代的问题：\textbf{如何根据有限证据推断应当修复什么，以及如何验证生成的补丁确实满足修复要求}~\cite{sokavr,smith2015cure}。

大语言模型（Large Language Model，LLM）为处理漏洞描述、源代码和修复经验提供了新的技术途径。进一步结合代码检索、文件编辑、编译执行与测试工具，模型能够在智能体框架中逐步完成定位、修改和反馈驱动的修复。PatchAgent 等工作已经将这些环节整合到交互式系统中，推动研究对象从给定代码片段上的补丁生成扩展至真实项目中的自动修复~\cite{patchagent}。然而，交互能力的增强并不自动等同于修复语义的正确理解：系统仍需判断告警位置与根因位置之间的关系、识别跨函数依赖，并区分补丁通过当前验证与恢复完整安全要求之间的差别。

Rust 等系统编程语言使这一问题具有更丰富的语义层次。RustBelt 从形式化角度说明，基于 \texttt{unsafe} 实现的安全抽象需要满足相应验证条件，才能作为安全语言的可靠扩展~\cite{jung2018rustbelt}；针对真实 Rust 程序的研究也揭示了内存与线程安全实践中的复杂问题~\cite{qin2020rust}。因此，Rust 漏洞修复不仅涉及局部条件检查，还可能涉及所有权、生命周期、类型与 trait 约束、异常清理和并发状态。基于上述背景，本节围绕传统修复方法、LLM 与智能体修复、补丁正确性评价及 Rust 修复进展展开，重点讨论从漏洞表现恢复根因语义、从候选修改建立可信修复证据所面临的挑战。

\subsubsection{研究现状与趋势}

\paragraph{（一）基于搜索、程序合成与语义约束的传统修复}

传统 APR 为漏洞自动修复提供了候选搜索、修复位置选择和测试反馈等基础机制。Le Goues 等提出的 \textbf{GenProg} 通过遗传编程对程序进行变换，以失败测试与功能测试约束候选程序的选择~\cite{genprog}。Nguyen 等提出的 \textbf{SemFix} 将通过给定测试的要求编码为约束，并结合符号执行、约束求解和程序合成搜索修复表达式~\cite{nguyen2013semfix}。二者分别代表基于搜索和基于语义约束的修复思路，但主要面向一般程序缺陷，不能直接作为已覆盖多类安全漏洞的证据；其生成结果的可靠性还取决于所使用的测试、约束及允许的程序变换。

面向安全缺陷，研究进一步尝试从漏洞执行中提取比单次测试结果更强的修复条件。Gao 等提出的 \textbf{ExtractFix} 从崩溃或漏洞触发中提取约束，通过程序分析将其向候选修复位置传播，再合成满足约束的修改~\cite{gao2021extractfix}。该方法不只要求原始输入不再崩溃，而是将崩溃所违反的条件转化为修复义务，探索对输入变化更具一般性的修复。Zhang 等提出的 \textbf{VulnFix} 则利用程序状态快照、状态变异与归纳推断学习区分安全和危险状态的不变式，并据此生成补丁~\cite{zhang2022vulnfix}。这些研究说明，安全条件建模和反例反馈并非 LLM 出现后才引入的思想，传统语义方法已经为约束驱动修复提供了基础。

不过，这类方法的分析与合成建立在具体的漏洞模型、状态表示和候选修复空间之上。对于缺乏通用崩溃判据的权限、协议和状态逻辑问题，或者依赖复杂库语义的跨过程缺陷，单一崩溃约束难以完整表达修复要求；大型项目的路径、别名与环境建模也会增加分析成本~\cite{sokavr,gao2021extractfix,zhang2022vulnfix}。因此，传统方法的局限不宜简单归结为"缺乏语义分析"，而应理解为：已有语义表示与求解机制如何适应更复杂的安全要求，以及如何在精度与可扩展性之间取得平衡。

\paragraph{（二）从历史补丁学习到大语言模型生成}

学习式 AVR 将漏洞代码到修复代码的映射作为建模对象，从历史修复数据中获得候选生成能力。Fu 等提出的 \textbf{VulRepair} 基于 T5 类编码器—解码器架构，结合预训练与子词表示学习漏洞修复模式~\cite{fu2022vulrepair}。Zhou 等提出的 \textbf{VulMaster} 进一步扩展输入范围和信息来源，缓解代码截断及修复上下文不足的问题~\cite{zhou2024vulmaster}。这类方法降低了逐类手工设计补丁模板的需求，但仍需考虑训练数据覆盖、输入信息完整性以及评价方式的影响；与参考补丁相似，并不直接等同于恢复了相应安全性质。

随着通用代码模型的发展，研究开始探索不进行任务专门训练的修复方式。Pearce 等系统考察了\textbf{零样本漏洞修复}，通过组织漏洞代码与提示信息引导模型生成候选修改，展示了预训练模型迁移至安全修复任务的可能性~\cite{pearce2023zeroshot}。与面向固定数据分布的训练方式相比，提示式方法便于接入不同模型和额外上下文；但模型输出仍需接受程序级检查，提示中缺失的调用关系、错误处理约定或边界条件，不会因为模型具有代码生成能力而自动得到可靠补全。

训练式方法也在从修复结果拟合转向显式推理能力学习。Wen 等提出的 \textbf{Vul-R2} 结合领域推理数据、监督微调与课程式奖励训练，增强模型对漏洞成因及修复过程的建模~\cite{wen2025vulr2}。需要注意，其训练中的可验证奖励包含问答与字符级匹配等信号，并不等同于对完整仓库执行安全性质验证。总体而言，训练式与提示式方法是互补路线，关键问题不是预设某一路线必然更优，而是判断所获得的推理和生成能力能否转化为可应用、可执行且安全的补丁。

\paragraph{（三）上下文增强、根因分析与经验引导修复}

直接向模型提供漏洞函数并要求修改，容易遗漏决定修复策略的跨函数关系与领域约定。针对这一问题，Nong 等提出的 \textbf{APPATCH} 通过静态依赖分析提取漏洞相关代码，结合根因分析、自适应示例选择和渐进式提示生成补丁~\cite{nong2025appatch}。其根因分析能够根据模型提出的信息需求逐步扩展相关函数切片；示例选择也不只依赖代码相似性，而是比较漏洞成因，以引导后续修复。该工作已经将跨过程上下文、根因推理和历史经验整合起来，并非仅在固定函数文本上进行一次性生成。

APPATCH 同时显示了这一技术路线的边界：其主要流程以漏洞相关位置及类别为分析起点，论文也进一步考察了与 CodeQL 告警结合的端到端设置，因此不能将其概括为完全不能处理检测器输出。其失败分析仍涉及根因判断错误、上下文不足和边界处理不完整等问题~\cite{nong2025appatch}。独立的 \textbf{InterPVD} 研究则明确区分漏洞触发语句与待修复语句，并指出二者可能位于不同函数~\cite{li2024interpvd}。这意味着"定位到漏洞表现"与"识别能够恢复安全条件的修改位置"需要分别评价。

从这些研究可以归纳出两个进一步的问题。其一，上下文扩展既要补齐关键定义、调用关系和状态依赖，又要避免将大量无关实现交给模型；仅扩大输入长度不能替代有针对性的证据组织。其二，若早期根因解释存在错误，以该解释选择示例和生成补丁可能继续强化同一假设。因而，根因分析不仅需要给出自然语言说明，还需要能够对照源码与执行证据检验其成立条件，并在证据冲突时修正解释。这是对上下文增强与因果示例选择流程的进一步研究要求，而不是否认已有方法已经开展根因分析~\cite{nong2025appatch,li2024interpvd}。

\paragraph{（四）面向真实仓库的交互式修复智能体}

与将 LLM 作为局部生成模块的流水线不同，修复智能体通过持续观察项目、选择工具和利用执行反馈组织分析过程。\textbf{SWE-agent} 研究智能体—计算机接口如何支持仓库导航、文件编辑和程序执行，表明工具接口设计会影响模型完成软件工程任务的方式~\cite{sweagent}。\textbf{RepairAgent} 则面向程序修复，使模型在信息收集、修复素材检索、候选修改和测试之间自主选择操作~\cite{repairagent}。这些工作主要提供通用软件工程或 APR 能力，不能仅凭通用基准上的任务完成率推断其在安全漏洞修复中的正确性。

\textbf{PatchAgent} 面向真实漏洞修复，将定位、补丁生成与验证整合到同一智能体中，并通过语言服务器及验证接口提供代码导航和反馈~\cite{patchagent}。这种设计使系统不必完全依赖预先固定的人工修复路径，能够根据构建失败、漏洞复现和回归测试调整候选补丁。不过，智能体能否持续执行修复，与反馈能否区分根因修复和表现抑制，是两个不同问题。多轮迭代可以提升对现有测试的适配程度，却不能单独保证测试未覆盖的行为也得到正确处理。

2026 年公开的 \textbf{Kumushi} 进一步将根因定位作为修复系统的核心，通过多样化动态定位与证据加权排序，为模型筛选更有诊断价值的程序位置；其评价同时结合自动验证与结构化专家审查，以区分通过 oracle 的补丁与更充分的根因修复~\cite{wang2026kumushi}。这一进展表明，研究重点已不仅是增加智能体的自主操作范围，而是改善"程序证据—根因判断—修复决策"之间的联系。与此同时，自动执行、人工审查和形式化证明具有不同的证据强度，不能因系统包含分析工具或多个角色，就将其结果视为已获得充分正确性保证。

\paragraph{（五）从可复现修复任务到独立补丁验证}

修复方法的发展依赖能够复现漏洞及验证修改的实验基础。\textbf{Vul4J} 提供带有漏洞证明测试的 Java 漏洞案例，支持在真实项目中研究自动修复~\cite{bui2022vul4j}；Wu 等通过 \textbf{VJBench} 及相关实验考察神经网络修复安全漏洞的能力，为区分一般缺陷与安全修复需求提供了数据基础~\cite{wu2023vjbench}。\textbf{SEC-bench} 将任务扩展至真实 C/C++ 仓库中的漏洞复现与修复，并通过智能体自动完成相当部分案例构建工作~\cite{secbench}。因此，现阶段的问题已经不只是缺少案例或依赖人工打包，而是如何在自动扩展基准的同时确保环境、漏洞证据和验证逻辑可靠。

补丁验证方面，Smith 等较早揭示了 APR 的测试过拟合问题~\cite{smith2015cure}。Yang 等提出的 \textbf{OPAD} 利用自动生成的测试输入与运行时安全判据，增强对过拟合补丁的识别~\cite{yang2017opad}。其意义在于，将用于产生补丁的有限测试之外的行为纳入检查；但增加输入数量主要扩展被执行的行为范围，能否发现错误仍取决于输入是否触及遗漏路径，以及 oracle 是否能够识别相应违规。针对权限策略、资源约束或接口契约等问题，仅观察崩溃并不足以判定修复质量。

2026 年的 \textbf{PVBench} 研究进一步引入开发者随官方修复新增的 \textbf{PoC\textsuperscript{+} 测试}，比较原有功能测试与 PoC 所形成的基础验证，以及包含修复后行为要求的增强验证~\cite{yu2026pvbench}。其结果显示，在所评估系统和案例上，基础验证接受的部分补丁仍会违反新增测试所表达的要求。这说明修复评价需要关注漏洞相关的预期行为，而非仅判断已知触发是否消失。另一方面，PoC\textsuperscript{+} 也可能包含实现策略、性能或工程惯例方面的要求，因此使用这类检查时，应区分漏洞必要安全条件与更广泛的补丁质量要求，避免把所有偏离官方实现的方案统一判为不安全。

同年公开的 \textbf{PatchBench} 则同时关注补丁记忆、表层修复和功能保持：通过漏洞移植与代码变换调整任务上下文，利用额外漏洞输入检查安全性，并在正常输入上对照参考修复后的程序输出及项目测试检查语义~\cite{shen2026patchbench}。该研究也指出，验证仍依赖有限输入，且真实修复时通常没有可直接用于比较的参考修复程序。由此可见，\textbf{更强的基准评估与修复过程可获得的反馈之间仍存在信息差}：评估者可以保留官方补丁和独立测试，而修复器必须在这些信息不可见时形成有效修改；二者的信息边界需要明确区分。

在上述背景下，本课题组前期构建的 \textbf{RRBench}~\cite{rrbench} 针对 Rust 生态缺乏大规模真实漏洞修复基准的问题，构建了覆盖 115 个案例、95 个仓库的仓库级修复基准，其两阶段评估设计与本节讨论的方向直接对应：公开阶段向修复系统提供漏洞触发 PoC 与基础功能测试，私有阶段使用对修复系统不可见的安全性质检查（security oracles）与合法行为检查，验证补丁是否真正恢复应有安全性质；对五款代表性修复智能体的评估表明，通过公开验证的提交中有 23.4\% 未能通过私有安全性质检查，且存在 14 个任何系统均无法正确修复的案例。这说明功能测试会系统性高估漏洞修复能力，补丁验证必须显式建模安全性质，也为研究内容二"生成—验证"协同机制的探索提供了直接的实证动机。

\paragraph{（六）Rust 修复及其语言相关验证机制}

Rust 修复研究已经覆盖多种任务，但不同工作的目标不能混为一谈。\textbf{RustAssistant} 利用编译器诊断与 LLM 迭代修复编译错误~\cite{deligiannis2025rustassistant}；\textbf{RUPAIR} 面向缓冲区溢出的检测与修复~\cite{hua2021rupair}；\textbf{RUSPATCH} 研究 Rust 应用的动态补丁部署，重点在于及时应用修复，而非从漏洞描述中自主合成补丁~\cite{ruspatch}。这些工作分别回答编译可接受性、特定缺陷修复和补丁应用问题，不能据此推导出对任意 Rust 安全漏洞的端到端修复能力。

在仓库级任务上，Xiang 等提出 \textbf{Rust-SWE-bench} 和 \textbf{RustForger}，通过自动测试环境设置及基于 Rust 元编程的动态跟踪，增强问题复现与程序分析~\cite{xiang2026rustswebench}。该工作面向一般 Rust 软件问题，其语言相关支持有助于仓库级修复，但基准任务完成并不自动证明漏洞安全要求已经恢复。\textbf{AkiraRust} 则面向未定义行为修复，使用有限状态机协调多智能体的快慢推理、运行时语义检查与回退~\cite{akirarust}。这表明执行语义已经进入 Rust 智能修复流程，但未定义行为检查与协议、权限、资源等安全要求的验证仍需采用不同机制。

Rust 的分析工具为此提供了互补支持。\textbf{Miri} 能在执行中识别多类未定义行为，但其文档明确说明，通过某些测试不能保证安全抽象对任意合法客户端均成立~\cite{miri}。\textbf{Loom} 通过模型化并发操作探索线程交错，但需要被测程序及相关组件使用其可观察的替代接口，未纳入模型的操作不在其分析范围内~\cite{loomtool}。Rust 语言参考还明确指出，\texttt{unsafe} 不豁免未定义行为约束，非法值、引用和别名关系等条件可能在危险状态形成时就已被违反~\cite{rustref}。因此，修复验证需要依据具体性质组合编译期客户端检查、执行检查、并发测试和合法行为对照，而不是寻找一个适用于所有 Rust 漏洞的统一崩溃判据。

综合上述研究，基于智能体的 AVR 呈现出三个相互关联的发展方向：从给定局部代码的补丁生成扩展至真实仓库中的交互式分析；从单纯学习修改模式转向结合上下文、根因和程序证据组织修复；从已有测试通过转向独立的安全行为与合法功能验证。三者的交汇点在于：模型不仅要生成修改，还需要明确修改依据、作用边界及其验证范围~\cite{patchagent}。

\subsubsection{研究挑战}

\paragraph{（一）不完备定位条件下，根因位置与修复上下文的可靠恢复}

InterPVD 揭示了触发位置与待修复位置之间的跨过程差异，APPATCH 和 Kumushi 则分别从静态上下文与动态证据出发支持根因分析~\cite{li2024interpvd,nong2025appatch,wang2026kumushi}。然而，对于检测器仅提供崩溃位置、污点终点或局部告警的情况，系统仍需要追踪状态的建立、传播和使用关系，确定真正应承担安全检查或状态维护的程序位置。同时，扩大上下文可能引入噪声，过度裁剪又可能丢失别名、辅助函数及异常路径。关键挑战是将不确定的定位线索转化为可验证的依赖证据，并以这些证据决定何时扩展、裁剪或重新定位，而不是默认上游告警已经解决修复定位问题。

\paragraph{（二）根因推理、历史经验与候选策略之间的语义一致性}

已有方法已通过因果示例、推理训练及证据排序改善补丁生成，但根因解释与目标程序之间仍可能存在偏差~\cite{wen2025vulr2,nong2025appatch,wang2026kumushi}。同一 CWE 类别只能提供缺陷类型线索，不能唯一确定被约束的对象、异常来源或正确干预方式。后续研究需要让候选解释接受数据流、控制关系、接口前置条件及执行观察的检验，并在证据不足时保留竞争性假设。历史案例也应作为待核对的经验，而非直接套用的模板。与之对应，候选补丁的多样性应体现不同的约束建立方式或修复位置，而不仅是对同一修改反复采样；增加候选数量是否改善正确修复覆盖，需要在相同预算下单独验证。

\paragraph{（三）跨接口、状态与语言机制的安全性质恢复}

RustBelt、真实 Rust 安全问题研究及语言参考表明，安全要求可能跨越类型接口、内部实现和多次调用，而不只对应一条局部语句~\cite{jung2018rustbelt,qin2020rust,rustref}。例如，限制一次危险访问并不必然修复产生非法引用的接口，改变异常返回也不必然修正形成错误状态的转移。因此，修复需要确定安全性质应当在何处建立、由哪些操作共同维护，以及哪些调用、清理或并发路径仍可能绕过约束。对于生命周期、trait 和接口契约问题，增强编译期约束也可能是有效修复方式；评价时必须区分危险客户端被有意拒绝与库本身或合法客户端意外构建失败，不能仅以运行时 PoC 的执行结果概括修复成败。

\paragraph{（四）部分规约下，独立验证与合法行为保持的可执行化}

OPAD、PVBench 和 PatchBench 已从额外输入、开发者新增测试和输出语义等角度增强验证，但这些机制仍受测试覆盖、参考信息及 oracle 适用范围限制~\cite{yang2017opad,yu2026pvbench,shen2026patchbench}。进一步的难点在于：从漏洞描述和程序证据中提取确实必要的安全条件，将其映射为能够执行的检查，并同时保留与危险行为相邻的合法用法。验证不仅要防止检查过弱而接受不完整修复，也要防止检查过强而拒绝不同于官方实现的正确方案。对于自动生成的验证逻辑，还需检查其前置条件、预期结果和程序映射，避免用模型生成的另一段解释为候选补丁自我背书。

这里还需要严格区分修复期反馈与最终评价依据。修复期可以利用公开测试和自行生成的反例迭代，但最终保留测试不应反复返回给修复器，否则其独立性会被削弱。对于 Rust，Miri、Loom 与编译期检查可以提供不同侧面的反例或支持证据，但均需要说明模型、输入和环境范围~\cite{miri,loomtool}。\textbf{通过有限验证表示满足已检查的要求，而不是证明程序在所有可能执行中均无漏洞。}这一边界同样意味着，核实某个修复是否存在的 PPT 不能单独替代补丁充分性与功能保持验证~\cite{smith2015cure}。

\paragraph{（五）真实任务、评价标准与资源成本的统一刻画}

不同修复研究在输入粒度、定位提示、运行环境和判定标准上存在明显差异：有的评价给定漏洞位置上的候选生成，有的要求系统自行分析仓库；有的采用参考补丁比较或人工审查，有的依赖可执行测试~\cite{sokavr,nong2025appatch,secbench,yu2026pvbench}。因此，需要在明确输入与信息权限的前提下，分别报告补丁生成、可应用构建、公开测试及独立验证的结果，并保留无补丁、环境失败和预算耗尽等情况。不能将不同证据标准下的"正确率"直接排序，也不能把基准任务的成功等同于已经完成生产环境中的安全修复。

此外，历史修复可能出现在训练或检索数据中，模型与官方补丁相似并不足以区分经验泛化和直接记忆；PatchBench 对任务上下文变换的研究进一步说明了控制这一风险的必要性~\cite{shen2026patchbench}。面向可复现研究，还应固定源版本、工具链、模型配置和检索边界，并在成本评价中纳入上下文获取、候选生成、测试与复杂验证的完整开销。最终需要回答的不是某个模型在最有利设置下能够生成多少补丁，而是在有限预算和不完备证据条件下，系统能够多可靠地完成根因定位、生成有效修改并提供独立验证证据。

综上，基于智能体的漏洞自动修复已经具备从局部生成走向交互式仓库分析的技术基础，但可信修复仍依赖两个相互衔接的能力：一是从不完备定位与上下文中形成有证据支持的根因判断，并据此探索适当的修复策略；二是把漏洞相关安全要求转化为独立、可执行且兼顾合法行为的验证条件。本课题研究内容二围绕这一联系展开，重点研究根因推理、候选修复与安全性质验证之间的协同机制，以提升复杂真实漏洞修复的可靠性与可评价性。

\subsection{基于智能体的漏洞修复验证}

\subsubsection{研究背景}

漏洞修复是软件安全保障的重要环节，其目标不仅是消除已知漏洞的触发条件，还需要在维持合法功能的前提下阻断危险行为。围绕这一目标，自动漏洞修复（Automated Vulnerability Repair，AVR）研究形成了漏洞分析、补丁生成与补丁验证三个主要环节，采用程序分析、约束求解、程序变换和学习等技术降低对人工修复的依赖~\cite{sokavr}。随着大语言模型的发展，修复过程进一步由单一补丁生成扩展为交互式任务执行。例如，PatchAgent 将缺陷定位、补丁生成与验证组织到同一智能体中，利用语言服务器和验证工具支持多轮修复~\cite{patchagent}。这些研究主要回答"如何产生有效修复"，但软件系统的实际安全状态还取决于修复能否被正确集成并传播至最终部署的程序。

从软件供应链的角度看，漏洞修复远超上游提交补丁这一局部开发活动，它贯穿上游项目、下游分支、发行版、设备厂商与最终用户，构成一个持续维护过程。Li 和 Paxson 对安全补丁的研究表明，补丁开发需要同时考虑修复完整性与功能回归，安全修复也可能仅涉及少量代码变化~\cite{li2017largescale}。针对修复在相关仓库之间的传播，Machiry 等提出 SPIDER，通过分析补丁的行为影响，识别适合向相关项目传播的修复~\cite{machiry2020spider}。Zhang 等对 Android 内核补丁生态的研究进一步发现，修复在多层供应链中的传播可能存在较长延迟~\cite{zhang2021android}。因此，"补丁已经公开""下游已经合入"与"部署产物实际包含修复"是三个需要分别确认的状态。

版本信息和更新公告能够辅助判断修复状态，但不能在所有场景下替代程序级证据。一方面，下游维护者可能将安全修复回移植至较早版本，同时保留原有上游版本号。Red Hat 的安全回移植实践明确指出，仅依据上游版本号开展漏洞扫描可能产生误判~\cite{redhatbackporting}。另一方面，闭源产品、固件和第三方二进制组件往往无法提供与部署产物完全对应的源码，导致使用者难以直接检查安全修复是否进入最终程序。由此，针对实际二进制开展独立的修复状态核查，成为连接补丁发布与安全维护的重要技术需求~\cite{fiber}。

补丁存在性测试（Patch Presence Testing，PPT）正是针对这一需求提出的：给定已知漏洞的修复信息，以及补丁状态未知的目标程序，判断目标中是否包含对应修复的实现证据。不同方法使用的参考材料有所区别，包括源码补丁、修复前后源码、参考二进制或中间表示等~\cite{fiber,binxray}。PPT 与通用二进制相似性检测的区别在于，后者主要寻找功能或结构相近的代码，前者则需要在高度相似的两个实现之间识别安全相关差异。一个边界条件、错误分支或参数处理的变化，可能不足以明显改变函数整体相似度，却足以改变其漏洞状态~\cite{binxray}。

需要区分的是，\textbf{PPT 核查"特定修复是否体现在目标程序中"，而补丁正确性验证考察"修复是否充分消除目标缺陷并维持合法行为"。}即使检测到参考修复已经存在，也不能据此证明参考补丁本身完整，更不能推导出目标程序不存在其他漏洞~\cite{sokavr,li2017largescale,fiber}。因此，本节将 PPT 定位为漏洞修复生命周期中的部署侧验证环节，重点讨论其二进制证据提取、补丁语义匹配以及智能体化分析能力。

\subsubsection{研究现状与趋势}

\paragraph{（一）基于局部语法与结构差异的补丁存在性测试}

基于语法与结构差异的方法通过比较修复前后的程序变化构造补丁特征，再在目标中搜索相应特征。与直接比较整个函数相比，这类方法将分析范围收缩到补丁影响区域，以减少无关代码变化对判定的干扰~\cite{binxray,patchdiscovery}。

Xu 等提出的 BinXray 采用两阶段匹配：先通过函数签名检索疑似目标函数，再通过修复前后函数之间的基本块映射提取补丁签名，利用长度敏感的相似性比较区分漏洞版本与修复版本~\cite{binxray}。其重要特点是分析可以在二进制层面完成，不要求提供目标源码，并将函数检索与补丁匹配纳入同一流程。由此可见，传统方法已经开始关注目标定位，但函数级相似性只能提供候选范围，仍需进一步识别区分修复前后状态的补丁级证据。

PatchDiscovery 进一步关注补丁签名的选择性。该方法在规范化和简化的控制流图上匹配基本块，分别提取漏洞版本与修复版本的关键基本块，再在目标函数中比较其匹配情况~\cite{patchdiscovery}。与保留较多路径信息的签名相比，关键基本块能够减少无关指令和版本变化引入的噪声。这一路线的共同优势是特征直观、匹配流程相对明确，适合进行快速筛查；但局部结构特征仍需要在参考程序与目标程序之间保持可对应性。编译优化、基本块重组以及补丁之外的结构变化，都会使"同一修复"与"相似结构"之间的对应关系更加复杂~\cite{binxray,patchdiscovery}。

\paragraph{（二）基于符号执行与中间表示的语义分析}

为降低对指令形式的依赖，研究者开始利用符号执行和中间表示描述补丁引起的行为变化。Zhang 和 Qian 提出的 FIBER 结合源码变更分析与二进制符号执行，将选定的补丁位置转换为可匹配的二进制签名，并在匹配时结合局部控制流结构与语义公式~\cite{fiber}。该方法通过结合结构约束与语义约束，增强签名对二进制表示差异的适应能力。FIBER 的实际评估使用带有符号信息的内核镜像；对于符号不可用的情形，论文讨论了接入相似性检索进行候选函数定位，但未在该工作中完成这一集成。

Jiang 等提出的 PDiff 面向经过厂商定制的下游内核，通过锚点基本块限定补丁影响路径，再利用路径约束、内存状态和函数调用等信息生成语义摘要，比较目标与修复前后参考版本的接近程度~\cite{pdiff}。该设计将比较对象从表面指令扩展为补丁相关路径上的行为，能够减轻定制代码的干扰。不过，其目标函数定位使用内核中的 KALLSYMS 信息，路径枚举也需要采取循环展开等近似策略，因此仍存在输入条件和分析范围上的边界。

Yang 等提出的 Robin 则通过符号执行寻找通向补丁区域的路径，构造能够区分修复前后行为的函数输入，并比较相同输入下的执行特征~\cite{robin}。它强调使用补丁引起的行为差异校验相似函数，而非仅依赖整体相似度。与此同时，函数级欠约束执行中得到的路径和输入，不必然对应完整程序中的真实可达执行，这也说明"行为特征相似""补丁存在"与"漏洞可实际触发"需要分别讨论。

PS\textsuperscript{3} 通过符号模拟提取寄存器写入、内存存储、条件与函数调用等语义签名，利用参考与目标之间的签名比较降低编译选项变化的影响~\cite{ps3}。REACT 将源码编译得到的 LLVM IR 与目标二进制提升得到的 IR 纳入统一分析流程，通过语义特征提取、特征细化与排序，以及约束求解进行补丁判定~\cite{react}。这些方法拓展了 PPT 的语义比较能力，但分析精度仍取决于符号模拟、二进制提升和特征抽象的准确性。IR 恢复误差与不完整的执行模型可能影响后续匹配，因此，即使采用约束求解，比较结论仍受所提取特征与建模条件的限制。

\paragraph{（三）不同程序表示下的补丁检测与演化问题}

除原生二进制外，Java 和 Android 字节码也形成了相应研究路线。BScout 直接检查 Java 可执行文件中的完整补丁，突破了仅依赖少量局部签名的传统做法~\cite{bscout}；PPT4J 结合补丁的语义信息与字节码特征，建立源码补丁行和二进制实现之间的对应关系~\cite{ppt4j}；PHunter 面向 Android 应用中的混淆代码，通过由粗到细的定位与匹配提升补丁识别能力~\cite{phunter}。这些方法表明，分析对象保留的信息、代码变换方式以及定位粒度，会直接影响技术设计。字节码场景的能力不能不加区分地推广到经过优化和符号剥离的 C/C++ 原生二进制。

另一个需要独立考察的问题是安全补丁演化。Xie 等系统研究了补丁代码在后续软件版本中的持续变化，以及这些变化对已知漏洞分析工具的影响~\cite{patchevolution}。编译差异主要改变同一源码实现的机器级表达，而补丁演化可能直接改变修复实现，例如重组检查逻辑或调整相关代码。因此，能够抵抗编译选项变化，不等于能够识别经过长期重构的修复。由此可见，传统方法的关键问题不仅是生成更加稳定的签名，还包括厘清签名所代表的安全行为，以及该行为在不同实现中成立所需的条件。

\paragraph{（四）基于大语言模型的补丁语义分析}

大语言模型为连接高层修复信息与低层程序表示提供了新的技术途径。DeGPT 利用模型改善反编译结果的可理解性，LLM4Decompile 探索从二进制代码生成或改进类源码表示，ReSym 则结合模型与程序分析恢复剥离二进制中的变量和数据结构符号~\cite{degpt,llm4decompile,resym}。这些工作提供的是程序理解的辅助能力，并不直接回答特定安全补丁是否存在；它们的价值在于为跨表示的语义定位与比较提供了基础。

Li 等提出的 Lares 是直接面向 PPT 的 LLM 方法。它从修复前后源码提取补丁相关切片，补充数据依赖、控制依赖和宏信息，再由 LLM 将切片映射到目标函数的反编译伪代码；随后利用 Z3 检查可提取公式的等价性，并在不能直接判定时使用 LLM 比较语义~\cite{lares}。这种设计消除了为提取参考特征而重新编译历史版本的要求，将模型用于跨表示定位，而非仅作为最终二分类器。

Lares 的适用边界同样重要：\textbf{其方法明确假设潜在漏洞函数已由外部二进制相似性方法识别，所解决的定位主要限于目标函数内部的代码切片定位，尚未覆盖从完整二进制独立完成全部检索步骤的场景。}其失败分析还指出，重复检查与细微修改可能干扰判断~\cite{lares}。因此，"能够处理剥离后二进制的伪代码"与"不依赖任何定位提示完成端到端分析"并非同一能力；局部公式匹配也不能替代完整安全语义验证。

Xu 等提出的 PPTSP 采用另一种混合设计：先对语义等价的指令序列进行规范化，再提取补丁影响的关键控制路径以减少无关代码干扰，最后在结构特征相近时引入 LLM 判断语义等价性~\cite{pptsp}。与 Lares 将模型用于源码—伪代码映射不同，PPTSP 主要将模型用于传统特征处理之后的语义判别。两者体现出互补的结合方式：模型可以辅助定位候选证据，也可以处理已有特征难以区分的情况，但不同使用位置带来的收益与风险需要分别评估。

\paragraph{（五）从 LLM 增强流水线到交互式智能体分析}

将 LLM 加入某个分析阶段，与由智能体组织完整分析过程仍存在区别。SWE-agent 的研究表明，专门设计的智能体—计算机接口能够支持模型自主浏览仓库、执行程序并根据反馈调整操作~\cite{sweagent}。结合这一交互范式，可以进一步探索由智能体围绕补丁假设主动选择二进制分析工具，在候选检索、局部检查与证据补充之间迭代，避免始终沿固定顺序运行预定义模块。不过，通用代码任务与 PPT 在分析对象和证据要求上存在差异，这种交互能力能否转化为可靠的端到端补丁检测，仍需要专门研究。

这一变化还意味着评价对象需要从单次模型输出扩展为完整分析系统。模型、工具接口、上下文组织、停止条件与重试机制都会影响执行过程；有关智能体评价的研究也指出，仅关注任务准确率而忽视资源使用和复现条件，可能对系统实用性产生误判~\cite{sweagent,aiagentsthatmatter}。对 PPT 而言，需要进一步回答：模型能否自行找到正确代码，工具返回的信息能否构成补丁证据，失败发生在哪个环节，以及正确结论需要付出多少分析成本。因而，现有研究呈现出的发展方向已超越简单地"以 LLM 替代传统程序分析"，更强调将程序分析提供的可检查证据与模型的语义推理、任务规划能力结合，并对整条流程进行验证。

围绕上述问题，本人前期已开展两项相互衔接的实证研究。其一，针对现有 PPT 研究评估数据集有限且不一致、失败根因缺乏系统理解的问题，构建了包含 561 个 CVE、覆盖 10 个广泛使用开源项目的高保真基准，对五款代表性工具进行系统评估，揭示了"高准确率、低决策成功率"的普遍现象，并将失败模式归纳为算法层局限与工程层缺陷两大类~\cite{zhang2026pptempirical}。其二，针对智能体能否胜任端到端 PPT 这一问题，系统评估了通用编码智能体在符号可见、剥离二进制、演化补丁与真实发行版包等场景下的能力边界，发现前沿智能体仅使用 binutils 即可达到与专用 SOTA 系统相当的性能，但补丁定位与演化补丁识别仍是主要瓶颈~\cite{zhang2026agenticppt}。这两项工作为研究内容三的提出提供了直接的实证基础。

\subsubsection{研究挑战}

\paragraph{（一）缺少定位提示时，补丁语义与目标代码之间的可靠对应}

现有方法在定位方面具有不同前提：FIBER 和 PDiff 的评估或实现利用目标符号信息，BinXray 包含函数级检索，Lares 则假设目标函数已由外部方法识别~\cite{fiber,binxray,pdiff,lares}。这些差异意味着，已知目标函数上的匹配效果不能直接代表完整二进制上的端到端效果。面向实际部署，分析不仅需要找到语义相关的函数，还必须确定具体检查作用于哪个对象、处于哪个程序位置，以及是否约束了对应危险操作。后续研究需要将函数检索、程序点定位与语义验证联系起来，并在候选存在歧义时继续收集证据，避免直接接受首个相似结果。

\paragraph{（二）参考实现变化时，修复语义的保持与可观察性}

语法签名、符号执行摘要和 IR 特征分别缓解了不同形式的程序差异，但它们仍基于一定的实现对应关系~\cite{binxray,patchdiscovery,pdiff,robin,ps3,react}。补丁演化研究则表明，原始修复代码会在后续版本中持续变化~\cite{patchevolution}。因此，需要进一步区分"原有补丁模式不再出现"与"修复效果不再存在"：前者可能只是实现调整，后者才涉及修复缺失。针对这一问题，应围绕检查对象、必要前提、错误传播和危险操作之间的关系组织验证，并按需扩展至调用上下文，避免仅停留在提高局部签名的匹配容忍度。同时，当给定分析范围内缺乏可区分证据时，应保留判断边界，不能以模型推测弥补不可观察信息。

\paragraph{（三）模型推理与程序证据之间的一致性}

Lares 和 PPTSP 将模型用于定位或语义判别，说明自然语言推理能够参与传统分析难以直接处理的环节~\cite{lares,pptsp}。但模型生成的解释、反编译输出以及求解器结论具有不同的证据性质：解释需要对应实际代码，伪代码依赖反编译恢复质量，求解结论则受公式建模范围约束~\cite{react,degpt,resym}。后续研究需要明确每个结论依赖哪些地址、语句、控制关系和数据关系，并针对竞争性解释寻找反证。还应区分未发现修复、漏洞相关功能不适用和分析未完成，避免将工具失败或检索无结果直接转换为漏洞状态判断。

\paragraph{（四）不同研究条件下，能力与失败机制的可比较性}

代表性 PPT 研究分别采用内核、通用 C/C++ 项目、Java 字节码或 Android 应用，并在参考材料、定位支持与编译条件上具有差异~\cite{fiber,binxray,patchdiscovery,pdiff,robin,ps3,react,bscout,ppt4j,phunter}。因此，论文中的单一准确率或 F1 值不能脱离实验设置进行直接排序。需要在统一条件下分别考察特征生成、目标定位、语义判定与完整流程，明确是否提供符号、候选函数和调试信息，并同时记录错误判定、无法判定、超时与工具异常。只有保留这些失败情形，才能进一步判断局限源于补丁语义不可区分、实现假设不满足，还是工程处理过程不完整，避免将所有失败混为模型或算法精度问题。

\paragraph{（五）智能体分析效果、工具配置与成本之间的权衡}

交互式分析扩大了系统可以探索的证据空间，也增加了工具调用、上下文维护和重复推理的成本。通用智能体研究已经揭示接口设计和资源评价的重要性~\cite{sweagent,aiagentsthatmatter}，但不能由此预设 PPT 中"更大模型""更多工具"或"更长上下文"必然带来更好效果。需要针对漏洞元数据、模型配置、工具能力与任务预算开展受控比较，研究如何优先使用低成本线索缩小范围，再针对关键证据缺口调用复杂分析，并同时评价正确性、完成率、时延和资源消耗。

综上，现有研究已经从补丁局部差异匹配发展到符号语义分析，并开始引入 LLM 进行跨表示定位与语义判别。然而，面向实际部署产物的可靠 PPT，仍需要贯通补丁理解、无充分提示的目标定位、演化修复识别和证据约束下的结论生成。对这些环节进行统一评估与失败机理分析，并探索程序分析和智能体推理之间的有效分工，是本课题研究内容三的主要切入点。

